% Modernised reading — fonds Alexandre Grothendieck, shelfmark 60, pages 1-69.
%
% This is an INTERPRETATION, not a transcription. It re-reads the pages in
% today's notation and vocabulary, states the hypotheses the manuscript leaves
% implicit, and drops the critical apparatus so that the mathematics reads
% straight through. Everything the brackets and underlines carried — a doubtful
% reading, a notation he used differently, a missing passage — moves into a
% footnote.
%
% It is held to being mathematically correct as it stands, not merely faithful
% to the page. Where the manuscript is loose or mistaken, the true statement is
% what appears here, and the footnote says what the page has. In any dispute of
% substance the transcription governs: whoever cites Grothendieck must cite the
% batch-NN.fr.tex files.
%
% Scope: the folder was read WHOLE before any of this was written — all four
% batches, pages 1-69 (thirty-four rectos in his hand, the versos being typed
% Bourbaki material skipped by the transcription) — and it is written as ONE
% file for the whole shelfmark.
%
% Notation fixed for the whole folder (spine section): F for Frobenius (his f,
% which he also uses for the structural morphism f : X -> Spec F_p — f is kept
% for the morphism only); r for the rank of a sheaf on a point (his n on p. 8,
% where n later is dim X); q = p^{n+rho} (q = p^{1+rho} for a curve);
% eps(E) := delta(E) q^{-chi(E)/2} throughout, and sgn(E) for the product of
% virtual signs of p. 22, which he also calls eps(E); m_+ , m_- for the
% multiplicities he writes eps_1, eps_{-1} (p. 16); on a curve, X° for the
% open set and u : X° -> X (his U and i), I_1 subset I_2 for his subgroups
% V subset U (pp. 38-46); L^natural_{E_eta} := L_{E_eta^natural}; delta_x,
% chi_x are taken after pushing forward to Spec F_p.
%
% Substantive departures from the pages, each footnoted where it occurs:
%   - p. 12: in (2) the exponent of (-p^{(n+rho)/2}/t) is +chi(E), not
%     -chi(E); (3), which he derives from it, is right; also the page's
%     « 2n Σ_{i<n} » is correct (the transcription's line has n);
%   - p. 24: D R_infty f(E) = R_infty f(DE)[-1] (his « 1 ou -1 ! »), and
%     [2n-1] rather than [1] once DE = E(n+rho)[2n] is substituted; the
%     parity, which is all the L-function sees, is unchanged;
%   - p. 26: eps(E) = delta'(E)^{-1} q^{chi/2}, not delta'(E) q^{-chi/2};
%   - p. 30: q = p^{n+rho}, not p^{n rho};
%   - p. 32 and p. 48: the shifts [2] are dropped on the page (even, harmless);
%   - p. 52: the local factors enter as lambda_x(qt) (lambda_x(pt) before the
%     weight substitution), the page writes lambda_x(E) with no argument;
%   - pp. 48-58: mu_x(E) is in fact zero for every E (proof given in the
%     body, ours); so D(E^natural) = (E-check)^natural(1) in the Grothendieck
%     group, lambda_x = 1, and the « ?? » of p. 52 is answered: yes;
%   - p. 22: « eps_i = +1 for i odd » and « i beta_i / 2 integer » need an
%     alternating Frobenius-compatible pairing on the odd-weight pieces (a
%     polarisation), not merely Weil's conjectures; stated under it;
%   - p. 8 and throughout: eps(E) = ±1 needs rational (or real)
%     characteristic polynomials; stated as a hypothesis.
%
% Announced and never established: the term of « type <= tau + (n-1) » of
% p. 20 (« type » is never defined); problems a), b), c) of p. 54 (only c) is
% reduced to b), p. 58); the identification of p. 60, left on a dash; the
% general case of pp. 62-68, abandoned at n = 1.
%
% Pass: Opus 5.5 (claude-opus-5-5), 2026-10-04 — first pass, unchecked
% against the pages by a human.
\documentclass[11pt,a4paper]{article}
\input{../preamble/grothendieck}

\folder{60}
\pages{1}{69}
\dating{[à partir de 1962- à partir de 1971]}
\watermark{Édition de démonstration\\interprétation personnelle de l'œuvre}
\foldertitle{Fonctions L / Équation fonctionnelle des fonctions L (cas géométrique) : notes manuscrites (s.d.) — lecture modernisée du dossier entier}

\begin{document}

\section*{Résumé}

\begin{resume}
La chemise porte, de sa main, « Notes sur l'équation fonctionnelle ». Ce sont
trente-quatre feuillets de brouillon, écrits au dos de tirages dactylographiés
de Bourbaki, qui cherchent à comprendre d'où vient la symétrie d'une
\emph{fonction~$L$} et ce qu'elle coûte quand les choses ne sont pas parfaites.

Une fonction $L$ est une série qui compte. On se donne une variété $X$ définie
sur un corps fini $\mathbb{F}_p$ --- un système d'équations polynomiales à
coefficients modulo $p$ --- et, en chacun de ses points, un peu d'algèbre
linéaire (un « faisceau » $E$) ; on fabrique une série $L_E(t)$ qui enregistre
cette algèbre linéaire point par point. Depuis Weil, on sait que ces séries
sont des fractions rationnelles et l'on s'attend à ce qu'elles soient
\emph{symétriques} : changer $t$ en $1/(qt)$, pour un certain nombre $q$,
redonne la même fonction, à un facteur simple près. C'est l'\emph{équation
fonctionnelle}, et le facteur, un signe et une puissance, est ce qu'on veut
connaître.

D'où vient la symétrie ? L'image est celle d'un miroir. La série $L_E$ se lit
sur la cohomologie de $X$ à coefficients dans $E$, et la dualité de Poincaré
dit que la cohomologie en degré $i$ est le reflet, retourné, de celle en
degré $2n-i$. Retourner un espace vectoriel, c'est inverser les valeurs propres
de ce qui agit dessus ; inverser les valeurs propres, c'est changer $t$ en
$1/t$ dans un polynôme caractéristique. Le premier tiers du dossier met cette
image en équations, avec soin : ce que devient le polynôme caractéristique
d'un dual (pages 2 à 6), ce que donne le miroir quand $X$ est propre et lisse
(pages 8 à 16), et le signe qui reste, qu'on sait déterminer quand
$n + \rho$ est impair et seulement en partie sinon.

Puis le miroir se fêle. Si $X$ n'est pas complète, il lui manque un bord, et
la symétrie ne tient plus qu'à un facteur près, porté par ce bord --- la
« cohomologie à l'infini » ; les pages 18 à 30 calculent ce facteur et
montrent qu'il vérifie à son tour sa propre équation fonctionnelle. Si $X$ est
une courbe et que $E$ n'est donné qu'au point générique, ce sont les points de
\emph{ramification} qui gâtent tout, et pire : la fonction $L$ naïve n'est même
pas additive. Les pages 32 à 58 inventent, pour réparer cela, une opération de
\emph{moyenne} sur le groupe d'inertie, notée $E \mapsto E^{\natural}$, qui
rend la fonction $L$ additive et locale, et calculent sa dualité. Cette lecture
montre en note que les termes correctifs qu'il y trouve sont en fait nuls, ce
qui répond par l'affirmative à la question qu'il laisse avec deux points
d'interrogation page 52. Enfin, pages 62 à 68, une réflexion abstraite :
étant donnée une équation $L(1/qt) = A(t)L(t)$, peut-on corriger $L$ par un
facteur « de niveau inférieur » pour que $A$ devienne une constante ?

Le ton est celui du travail en train de se faire : « $1$ ou $-1$ ! » en marge
d'un décalage, « Ceci n'est pas encore satisfaisant », « il vaut sans doute
mieux ne pas faire ce changement », un « ?? » encadré d'une flèche, et, page 28,
une vérification de cohérence faite pour elle-même --- substituer deux fois, et
retrouver l'équation du bord. Les conjectures de Weil y sont des hypothèses,
« en admettant la résolution des singularités … et les conjectures de Weil ».

Les noms à chercher ensuite : formule des traces de Grothendieck--Lefschetz,
dualité de Verdier, conjectures de Weil et leur démonstration par Deligne,
signe de l'équation fonctionnelle, formule de Grothendieck--Ogg--Chafarevitch,
théorème de monodromie $\ell$-adique, représentations de Weil--Deligne,
constantes locales.

\keywords{L-function of an l-adic sheaf, functional equation, Grothendieck--Lefschetz trace formula, Verdier duality, Poincaré duality, Weil conjectures, root number, big Witt vectors, lambda-ring, boundary cohomology, cohomology at infinity, Grothendieck--Ogg--Shafarevich formula, l-adic monodromy theorem, quasi-unipotent inertia, Weil--Deligne representation, local duality, Hilbert's Theorem 90}
\end{resume}

\section*{Le fil du dossier, et les conventions}

\pagerange{1}{1}
\subsection*{Les stations}

\begin{itemize}
  \item \textbf{Page 1} --- la chemise, « Notes sur l'équation fonctionnelle ».
  \item \textbf{Pages 2 à 6} --- algèbre linéaire : $P_F(t) = \det(1 - Ft)$,
    son comportement sous $\otimes$, $\lambda^i$, $\mathrm{Sym}^i$ et sous la
    contragrédiente.
  \item \textbf{Pages 8 à 10} --- un faisceau sur un point : poids,
    autodualité $\check E \simeq E(\rho)$, équation fonctionnelle de $L_E$.
  \item \textbf{Pages 10 à 16} --- $X$ propre : dualité globale, l'équation
    (1) en général, puis (2), (3) et le signe $\varepsilon'(E)$ quand $X$ est
    propre et lisse et $E$ autodual.
  \item \textbf{Pages 18 à 30} --- $X$ lisse non nécessairement complète : le
    facteur $A(t)$ venu de l'infini, la décomposition par poids (page 22),
    l'équation fonctionnelle du terme à l'infini (pages 24 et 26), une liste
    d'invariants arithmétiques (page 26), une vérification (pages 28 et 30).
  \item \textbf{Pages 32 à 36} --- une courbe, un faisceau au point
    générique ; la fonction $L^{*}$ et son défaut d'additivité.
  \item \textbf{Pages 38 à 46} --- « Suite de la remarque » : la moyenne
    $E \mapsto E^{\natural}$ sur l'inertie, sa formule de trace, et le
    faisceau virtuel $E_\eta^{\natural}$.
  \item \textbf{Pages 48 à 58} --- la dualité de $E_\eta^{\natural}$, les
    termes locaux $\alpha_x$, $\mu_x$, l'équation fonctionnelle qui en résulte,
    trois problèmes, et la réduction du troisième au deuxième.
  \item \textbf{Page 60} --- « Cohomologie à l'infini » : la suite exacte longue
    et deux triangles.
  \item \textbf{Pages 62 à 68} --- simplifier une équation fonctionnelle par un
    facteur de « niveau » inférieur ; le cas $n = 1$, commencé.
\end{itemize}

Le dossier est un seul mouvement, du plus simple au plus abîmé : un point,
puis $X$ propre et lisse, puis $X$ ouverte, puis une courbe ramifiée. La
page 60 revient en arrière sur la cohomologie à l'infini des pages 18 à 26, et
les pages 62 à 68 reprennent l'équation des pages 52 à 54 sous forme abstraite.
Aucune page ne porte de numérotation de sa main ; l'ordre est celui des
archivistes.\footnote{Les versos (pages impaires 3 à 69) sont des tirages de
rédactions de Bourbaki, sans rien de sa main ; la transcription les saute
sans les décrire et cette lecture fait de même. Ils donnent une borne
inférieure à la date d'écriture, que l'inventaire de Montpellier fixe à
« [à partir de 1962 -- à partir de 1971] ».}

\subsection*{Les conventions, valables pour tout le dossier}

$p$ est premier, $\ell \neq p$, et les faisceaux sont des
$\mathbb{Q}_\ell$-faisceaux constructibles (ou des complexes de tels
faisceaux). On note $F$ le Frobenius \emph{géométrique}, qui agit par $p^{-1}$
sur $\mathbb{Q}_\ell(1)$ : c'est la convention de la page 8, où
$L_{\mathbb{Q}_\ell(\rho)}(t) = (1 - p^{-\rho}t)^{-1}$.\footnote{La page note
le Frobenius $f$, et note aussi $f$ le morphisme structural
$X \to \operatorname{Spec}\mathbb{F}_p$ ; on réserve $f$ au morphisme
($Rf_{*}$, $Rf_{!}$) et l'on écrit $F$ pour le Frobenius.} Pour un
endomorphisme $F$ d'un espace de dimension finie,
$P_F(t) = \det(1 - Ft)$ ; pour un complexe $K$ sur
$\operatorname{Spec}\mathbb{F}_p$,
\[
  \begin{gathered}
  L_K(t) = \prod_i \det\bigl(1 - Ft \mid H^i(K)\bigr)^{(-1)^{i+1}},\\
  \chi(K) = \sum_i (-1)^i \dim H^i(K),\qquad
  \delta(K) = \prod_i \bigl(\det F \mid H^i(K)\bigr)^{(-1)^i}.
  \end{gathered}
\]
Pour $f : X \to \operatorname{Spec}\mathbb{F}_p$ et $E$ sur $X$, on pose
$L_E = L_{Rf_{!}E}$, $\chi(E) = \chi(Rf_{!}E)$, $\delta(E) = \delta(Rf_{!}E)$
: par la formule des traces de Grothendieck, $L_E$ est le produit eulérien
$\prod_x \det(1 - F_x t^{\deg x} \mid E_{\bar x})^{-1}$.\footnote{Le dossier
emploie $L_E$ sans le définir, et passe librement du produit eulérien à
l'expression cohomologique : c'est la formule des traces de Grothendieck
(exposé Bourbaki n° 279, 1964 ; SGA~5 et SGA~4½). Le nom de « formule des
traces de Grothendieck--Lefschetz » est postérieur.} La dualité $D$ est celle
de Verdier, $DK = R\mathcal{H}om(K, Rf^{!}\mathbb{Q}_\ell)$ ; sur un point,
c'est le dual linéaire. $\check E$ est la contragrédiente (il écrit $E$ surmonté
d'un háček).

$n = \dim X$ ; le rang d'un faisceau sur un point s'écrit $r$.\footnote{Page 8,
la page écrit $n$ pour le rang ; dès la page 10, $n$ est la dimension de $X$.
On sépare les deux.} Quand $E$ est de poids $\rho$, on pose
\[
  \begin{gathered}
  q = p^{\,n+\rho} \quad (\text{donc } q = p^{1+\rho} \text{ pour une courbe}),\\
  \varepsilon(E) = \delta(E)\, q^{-\chi(E)/2},
  \qquad \varepsilon'(E) = (-1)^{\chi(E)}\varepsilon(E),
  \end{gathered}
\]
et l'on garde la racine carrée positive de $q$.\footnote{La lettre
$\varepsilon$ sert dans le dossier à trois choses : la dimension
$\varepsilon(f)$ (page 6), la caractéristique d'Euler d'un complexe (pages 20
et 26), et un « signe » qui n'a pas toujours la même définition --- le signe
de $\det F$ sur $H^n$ (pages 8 à 16), le produit des signes de déterminants
virtuels (page 22), le quotient $\delta(E)q^{-\chi(E)/2}$ (pages 26 à 30). Les
trois dernières coïncident quand $X$ est propre et lisse et $E$ autodual ; en
général non (voir la note de la page 22). On écrit $r$ et $\chi$ pour les deux
premières, $\varepsilon(E)$ pour la dernière seulement, et
$\operatorname{sgn}(E)$ pour celle de la page 22.}

Une hypothèse court sous tout le dossier sans être écrite : pour que
$\varepsilon(E)$ soit un \emph{signe}, il faut que les polynômes
caractéristiques en jeu soient à coefficients rationnels (ou au moins réels) ---
ce que la page 8 appelle un \emph{motif}. Pour un $\mathbb{Q}_\ell$-faisceau
quelconque, pur de poids $\rho$, $\varepsilon(E)$ est seulement un nombre
algébrique dont tous les conjugués complexes sont de module~$1$. On le dit
une fois ici ; partout où l'on écrit $\varepsilon = \pm 1$, c'est sous cette
hypothèse.

\pagerange{2}{6}
\section*{1. Le polynôme caractéristique (pages 2 à 6)}

Soit $F$ un endomorphisme d'un espace $V$ de dimension $r$ sur un corps $k$,
de valeurs propres $\lambda_1, \dots, \lambda_r$ dans une clôture algébrique.
Alors
\[
  P_F(t) = \det(1 - Ft) = \prod_i (1 - \lambda_i t)
         = \sum_i (-1)^i \operatorname{Tr}\bigl(\Lambda^i F\bigr)\, t^i,
  \qquad
  \frac{1}{P_F(t)} = \sum_{\nu \ge 0} \operatorname{Tr}\bigl(\mathrm{Sym}^\nu F\bigr)\, t^\nu ,
\]
et les racines de $P_F$ sont les inverses des valeurs propres non nulles. $P_F$
est multiplicatif dans les suites exactes : c'est un invariant « additif » de
$V$, à valeurs dans le groupe multiplicatif $1 + t\,k[[t]]$.\footnote{En bas de
la page 2, un petit schéma de trois rangées d'indices ($\rho - 1, \rho,
\rho + 1$ ; $-2, -1, 0$ ; …) reliées par des ovales ; la page ne dit pas à
quoi il sert, et cette lecture ne l'interprète pas.}

Ce groupe porte une seconde loi, $*$, déterminée par
$(1 - at) * (1 - bt) = 1 - abt$ et la distributivité par rapport au produit
ordinaire : c'est l'anneau des \emph{grands vecteurs de Witt}, dans lequel
l'« addition » est le produit des séries.\footnote{Le nom n'est pas sur la
page ; la structure, oui, avec « la » $\lambda$-structure qui suit.} Alors,
avec $L_F = 1/P_F$,
\[
  P_{F \otimes F'} = P_F * P_{F'}, \qquad
  L_{F \otimes F'} = \bigl(L_F * L_{F'}\bigr)^{-1}, \qquad
  L_{F_1 \otimes \cdots \otimes F_\nu} = \bigl(L_{F_1} * \cdots * L_{F_\nu}\bigr)^{(-1)^{\nu+1}},
\]
la dernière formule résultant de ce que $L_F$ est l'opposé de $P_F$ pour
l'addition de Witt.\footnote{Le produit de la dernière formule est tracé
comme un $\prod$ ; c'est le produit $*$, seul à donner une formule juste.}
Pour la $\lambda$-structure, où $\lambda^j$ envoie $\prod(1 - \lambda_i t)$ sur
$\prod_{i_1 < \cdots < i_j}(1 - \lambda_{i_1}\cdots\lambda_{i_j}t)$, on a
$P_{\Lambda^j F} = \lambda^j(P_F)$, et la relation $\lambda^j(-x) =
(-1)^j\sigma^j(x)$, valable dans tout $\lambda$-anneau, donne
\[
  L_{\Lambda^j F}(t) = \bigl(\sigma^j L_F(t)\bigr)^{(-1)^{j+1}}, \qquad
  L_{\mathrm{Sym}^j F}(t) = \bigl(\lambda^j L_F(t)\bigr)^{(-1)^{j+1}} .
\]
\footnote{La page passe par une identité de séries génératrices en $s$ dont
une ligne est biffée et un facteur $(-1)^i$ surchargé ; on donne le résultat
encadré, qu'on a vérifié directement.}

Si $F$ est un automorphisme,
$P_{\check F}(t) = \prod(1 - \lambda_i^{-1}t) = (-t)^r(\det F)^{-1}P_F(t^{-1})$,
c'est-à-dire
\[
  P_{\check F}(t) = (-t)^{r}\,(\det F)^{-1}\,P_F(t^{-1}), \qquad
  L_{\check F}(t) = (-t)^{-r}\,\det F\; L_F(t^{-1}) .
\]
\footnote{Le signe $-$ de l'exposant de la seconde formule est douteux sur la
page ; il est juste. L'exposant y est noté $\varepsilon(f)$, la dimension.}
Il insiste : $t \mapsto t^{-1}$ n'a de sens que sur les séries qui viennent de
polynômes (de terme constant $1$, et de terme dominant inversible si $k$ est
un anneau) --- sur les fractions rationnelles, pas sur $1 + t\,k[[t]]$.

\pagerange{8}{10}
\section*{2. Un faisceau sur un point (pages 8 et 10)}

Soit $E$ un $\mathbb{Q}_\ell$-faisceau sur $\operatorname{Spec}\mathbb{F}_p$,
c'est-à-dire un espace de dimension $r$ muni de $F$. La section précédente donne
$L_{\check E}(t) = (-t)^{-r}\det F_E\,L_E(t^{-1})$. Supposons $E$ \emph{pur de
poids $\rho$} --- toutes les valeurs propres de $F$ de module $p^{\rho/2}$ dans
tout plongement complexe --- et à polynôme caractéristique rationnel. Alors
$\det F_E$ est un rationnel de module $p^{r\rho/2}$, donc
\[
  \det F_E = \varepsilon(E)\, p^{r\rho/2}, \qquad \varepsilon(E) = \pm 1,
\]
et $r\rho$ est pair : si $\rho$ est impair, $r$ est pair.\footnote{La
page dit « s'il s'agit d'un \emph{motif} satisfaisant aux conj.\ de Weil » ;
c'est la rationalité du polynôme caractéristique qui donne $\varepsilon = \pm 1$
et la parité, le poids seul ne suffit pas. Les mots qui suivent « dépendant du
motif $\Lambda^r E$ » ne se lisent pas.}

D'autre part $\check E$ et $E(\rho)$ ont les mêmes valeurs propres : celles de
$\check E$ sont les $\lambda_i^{-1} = \bar\lambda_i\,p^{-\rho}$, celles de
$E(\rho)$ les $\lambda_i p^{-\rho}$, et le multiensemble des $\lambda_i$ est
stable par conjugaison dès que $P_E$ est réel. On écrit
$\check E \simeq E(\rho)$ ; au niveau des fonctions $L$, c'est
exactement ce qui sert.\footnote{La page écrit $\check E \simeq E(\rho)$ « par
dualité ». Comme isomorphisme de faisceaux, cela demande un accouplement
$E \otimes E \to \mathbb{Q}_\ell(-\rho)$ parfait et compatible à $F$ ; pour les
fonctions $L$, l'égalité des polynômes caractéristiques suffit, et c'est elle
qu'on vient d'établir.} Ainsi
$L_{\check E}(t) = L_E(p^{-\rho}t)$,\footnote{La page le calcule par la loi $*$ :
$L_{E(\rho)} = (L_E * L_{\mathbb{Q}_\ell(\rho)})^{-1} = L_E * (1 - p^{-\rho}t)$.}
et en comparant,
\[
  L_E\Bigl(\frac{1}{p^{\rho}t}\Bigr) = \det F_E\,(-t)^{r}\,L_E(t)
  = \varepsilon(E)\bigl(-p^{\rho/2}t\bigr)^{r} L_E(t) .
\]
C'est le modèle de tout le dossier : la dualité donne $L_{\check E}$ en fonction
de $L_E(1/t)$, le poids donne $L_{\check E}$ en fonction de $L_E(t/p^\rho)$, et
l'équation fonctionnelle est leur rencontre.

\pagerange{10}{16}
\section*{3. Variété propre : l'équation fonctionnelle (pages 10 à 16)}

\subsection*{L'équation générale}

Soit $X$ propre sur $\mathbb{F}_p$ et $E$ constructible sur $X$. Le théorème de
dualité globale donne $H^i(\bar X, DE) \simeq H^{-i}(\bar X, E)^{\vee}$,
c'est-à-dire $Rf_{*}(DE) = D(Rf_{*}E)$. Appliquant la formule du dual de la
section~1 à chaque $H^i$ :
\[
  \text{(1)}\qquad
  L_{DE}(t) = (-t)^{-\chi(E)}\,\delta(E)\,L_E(t^{-1}),
\]
avec $\chi(E) = \sum (-1)^i b_i$, $b_i = \dim H^i(\bar X, E)$, \emph{invariant
géométrique}, et $\delta(E) = \prod_i (\det F \mid H^i(\bar X, E))^{(-1)^i}$,
\emph{invariant arithmétique} --- les deux mots sont dans la marge. Si de plus
$X$ est propre et lisse et $E$ lisse et pur de poids $\rho$, chaque
$H^i(\bar X, E)$ est pur de poids $i + \rho$, d'où
\[
  \delta(E) = \varepsilon(E)\, p^{\sum_i (-1)^i \frac{i+\rho}{2} b_i} .
\]
\footnote{Le numérateur $i + \rho$ est surchargé sur la page ; c'est bien le
poids de $H^i$. La pureté de $H^i(\bar X, E)$ est ce que la marge attribue
aux « conjectures standard » ; elle est démontrée depuis Deligne, \emph{La
conjecture de Weil}~I (1974) pour $E = \mathbb{Q}_\ell$ et II (1980) en
général. Le dossier la suppose, sans la tenir pour acquise. La page écrit tantôt
$E$, tantôt $F$, pour le même faisceau.}

\subsection*{Le cas favorable}

Supposons maintenant $X$ propre, lisse, purement de dimension $n$, $E$ lisse,
pur de poids $\rho$, et $\check E \simeq E(\rho)$ \emph{globalement sur $X$}.
La dualité de Poincaré donne $b_i = b_{2n-i}$, et
\[
  \sum_{i=0}^{2n} (-1)^i i\,b_i
  = 2n\sum_{i=0}^{n-1}(-1)^i b_i + (-1)^n n\,b_n = n\,\chi(E),
  \qquad\text{d'où}\qquad
  \delta(E) = \varepsilon(E)\,q^{\chi(E)/2} .
\]
\footnote{La page écrit bien $2n\sum_{i=0}^{n-1}$, et c'est juste (les termes
$i$ et $2n - i$ donnent ensemble $(-1)^i\,2n\,b_i$) ; la ligne de la
transcription, avec $n$, est à corriger.} Ici $\varepsilon(E) = \pm 1$, et
c'est le $\varepsilon(E) = \delta(E)q^{-\chi(E)/2}$ des conventions. L'équation (1)
devient
\[
  \text{(2)}\qquad
  L_{DE}(t) = \varepsilon(E)\Bigl(-\frac{q^{1/2}}{t}\Bigr)^{\chi(E)} L_E(t^{-1}) .
\]
\footnote{La page écrit l'exposant $-\chi(E)$. C'est un lapsus : de (1),
$(-t)^{-\chi}\varepsilon\,q^{\chi/2} = \varepsilon(-q^{1/2}/t)^{\chi}$, et
c'est avec $+\chi(E)$ que (2) donne (3), que la page écrit correctement.}
Or $DE \simeq \check E(n)[2n] \simeq E(n + \rho)[2n]$, d'où
$Rf_{*}(DE) \simeq Rf_{*}(E)(n + \rho)[2n]$ et, le décalage étant pair,
$L_{DE}(t) = L_E(t/q)$.\footnote{La page note en marge qu'il suffit que
$\check E \simeq E(\rho)$ hors d'un fermé de $X$ ; la fin de la phrase ne se
lit pas, et l'on n'en fait pas usage.} Portant dans (2) et changeant $t$ en
$t^{-1}$ :
\[
  \begin{gathered}
  \text{(3)}\qquad
  L_E\Bigl(\frac{1}{qt}\Bigr) = \varepsilon(E)\bigl(-q^{1/2}t\bigr)^{\chi(E)}L_E(t),\\
  \xi_E(t) = t^{\chi(E)/2}L_E(t), \qquad
  \xi_E\Bigl(\frac{1}{qt}\Bigr) = \varepsilon'(E)\,\xi_E(t),
  \end{gathered}
\]
avec $\varepsilon'(E) = (-1)^{\chi(E)}\varepsilon(E) = \pm 1$.

\subsection*{Le signe}

La dualité de Poincaré met $H^i(\bar X, E)$ et $H^{2n-i}(\bar X, E)(\rho)$ en
dualité à valeurs dans $\mathbb{Q}_\ell(-n)$ ; les déterminants de $F$ sur
$H^i$ et $H^{2n-i}$ ont donc même signe, les facteurs se compensent deux à
deux dans $\delta(E)$, et
\[
  \begin{gathered}
  \varepsilon(E) = \operatorname{sgn}\det\bigl(F \mid H^n(\bar X, E)\bigr),
  \qquad (-1)^{\chi(E)} = (-1)^{b_n},\\
  \varepsilon'(E) = (-1)^{b_n}\operatorname{sgn}\det\bigl(F \mid H^n(\bar X, E)\bigr).
  \end{gathered}
\]

\emph{Si $n + \rho$ est impair.} Supposons l'autodualité de $E$ donnée par un
accouplement $(-1)^\rho$-symétrique, comme pour la cohomologie de degré $\rho$
d'une variété.\footnote{Hypothèse implicite : la page invoque une « forme
alternée non dégénérée invariante » sans dire d'où vient son alternance.}
L'accouplement de Poincaré sur $H^n(\bar X, E)$ est alors alterné et non
dégénéré, et $F$ le multiplie par $q$ : $b_n$ est pair, et le déterminant d'une
similitude symplectique de multiplicateur $q$ vaut $q^{b_n/2} > 0$. Donc
$\varepsilon'(E) = 1$.\footnote{Il attribue l'argument, semble-t-il, à
Serre ; le nom est d'une lecture incertaine et l'on ne s'appuie pas dessus.}

\emph{Si $n + \rho$ est pair}, le signe n'est pas déterminé en général. Les
valeurs propres de $F$ sur $H^n$ non réelles vont par paires conjuguées de
produit $q > 0$ ; les réelles sont $\pm q^{1/2}$. Soit $m_{+}$, $m_{-}$ leurs
multiplicités. Alors $(-1)^{b_n} = (-1)^{m_{+} + m_{-}}$,
$\varepsilon(E) = (-1)^{m_{-}}$, et
\[
  \varepsilon'(E) = (-1)^{m_{+}} .
\]
\footnote{La page note $\varepsilon_1$, $\varepsilon_{-1}$ « la multiplicité de
la valeur propre $1$ ($-1$) de $f$ » : il faut l'entendre du Frobenius
normalisé $q^{-1/2}F$, entier ici puisque $n + \rho$ est pair. On renomme
$m_{\pm}$ pour ne pas confondre avec les signes $\varepsilon_i$ de la page 22.}

Hors du cas favorable, $DE$ ne s'exprime plus simplement en fonction de $E$ et
(1) relie deux fonctions, $L_{DE}$ et $L_E$, plutôt qu'une seule à elle-même.

\pagerange{18}{30}
\section*{4. Variété non complète : le terme à l'infini (pages 18 à 30)}

\subsection*{Le facteur $A(t)$}

Soit $X$ lisse, purement de dimension $n$, $E$ lisse sur $X$, pur de poids
$\rho$, avec $\check E \simeq E(\rho)$, et
$f : X \to \operatorname{Spec}\mathbb{F}_p$ compactifiable. La dualité
s'écrit maintenant $D\,Rf_{!}(E) = Rf_{*}(DE) \simeq Rf_{*}(E)(n + \rho)[2n]$.
Notons $Rf_{\infty}(E)$ le troisième sommet du triangle
$Rf_{!}E \to Rf_{*}E \to Rf_{\infty}E \to$, la « cohomologie à l'infini » que la
page 60 construira, et $L^{\infty}_E = L_{Rf_{\infty}E}$ ; dans le groupe de
Grothendieck, $[Rf_{*}E] = [Rf_{!}E] + [Rf_{\infty}E]$.\footnote{La lecture de
l'indice $\infty$ est celle de la page ; la page 24 écrit
$R_\infty f$ et $R^{*}_\infty f$ pour le même objet.} Les fonctions $L$ des deux
membres de la dualité donnent
$(-t)^{-\chi(E)}\delta(E)L_E(1/t) = L_E(t/q)\,L^{\infty}_E(t/q)$, c'est-à-dire
\[
  L_E\Bigl(\frac{1}{qt}\Bigr) = A(t)\,\delta(E)\,(-t)^{\chi(E)}\,L_E(t),
  \qquad A(t) = L^{\infty}_E\Bigl(\frac{1}{qt}\Bigr)^{-1},
\]
où $\chi$ et $\delta$ sont ceux de $Rf_{!}E$ ; si $X$ est complète,
$A = 1$.\footnote{Un encadré précise : « en admettant la résolution des
singularités … et les conjectures de Weil ».} Il ajoute que $L^{\infty}_E$ est
de poids $\le \rho + 2n$ et de « type » $\le \tau + (n - 1)$, $\tau$ étant le
type de $E$.\footnote{Le « type » n'est défini nulle part dans le dossier, et
l'on ne sait pas ce que cette borne affirme. La borne de poids est cohérente
avec ce qu'on sait aujourd'hui (les poids de $H^i$ d'une variété lisse de
dimension $n$ sont $\le 2n$) ; le dossier ne la démontre pas.}

Les deux membres de la dualité donnent aussi
$\chi(Rf_{*}E) = \chi(Rf_{!}E)$ --- le décalage $[2n]$ est pair --- et
$\delta(Rf_{!}E) = \delta(D\,Rf_{!}E)^{-1}$, d'où
\[
  \delta(E) = q^{\chi(E)}\,\delta(Rf_{*}E)^{-1},
  \qquad\text{et, si $X$ est complète,}\qquad \delta(E)^2 = q^{\chi(E)},
\]
« ce qui est bien compatible avec les conjectures de Weil sur les valeurs
absolues des valeurs propres ».\footnote{La page écrit
« $\chi(E) = \varepsilon(Rf_{!}E) = \varepsilon(D\,Rf_{!}E)$ » : son
$\varepsilon$ est ici la caractéristique d'Euler. L'égalité
$\chi(Rf_{*}E) = \chi(Rf_{!}E)$, qu'on tire ici de l'autodualité, est vraie
pour tout faisceau constructible : c'est le théorème de Laumon (1981).}

\pagerange{22}{22}
\subsection*{Décomposition par poids}

Page 22, pour un « motif » $E$ : décomposons la classe de $Rf_{!}E$ selon les
poids, $[Rf_{!}E] = \sum_i (-1)^i [h^i_{!}(E)]$ où $h^i_{!}(E)$ est la partie
virtuelle de poids $i$ et $\beta_i(E)$ son rang (« nombres de Betti
virtuels »). Alors $\chi(E) = \sum (-1)^i\beta_i(E)$ et
$\delta(h^i_{!}) = \varepsilon_i\,p^{i\beta_i/2}$, d'où
\[
  \delta(E) = \operatorname{sgn}(E)\,p^{\frac12\sum_i (-1)^i i\,\beta_i(E)},
  \qquad \operatorname{sgn}(E) = \prod_i \varepsilon_i(E) = \pm 1 .
\]
La page affirme que $i\beta_i/2$ est entier et $\varepsilon_i = +1$ pour $i$
impair. C'est vrai si chaque morceau de poids impair porte un accouplement
alterné non dégénéré que $F$ multiplie par $p^i$ --- ce que donne une
polarisation, et l'argument de la section~3 ; ce n'est pas une conséquence
des seules conjectures de Weil, ni de la rationalité.\footnote{Hypothèse
implicite, ajoutée. La page écrit « $\varepsilon(E)$ » pour ce produit de
signes ; on le note $\operatorname{sgn}(E)$, parce que la page 26 appelle
$\varepsilon(E)$ autre chose. La comparaison des deux est immédiate et n'est
pas sur la page :
$\varepsilon(E) = \delta(E)q^{-\chi(E)/2} = \operatorname{sgn}(E)\,
p^{\frac12(\sum_i (-1)^i i\beta_i - (n+\rho)\chi(E))}$ ; ils coïncident
exactement quand $\sum (-1)^i i\,\beta_i = (n + \rho)\chi(E)$, ce qui est le
calcul du cas favorable.}

\pagerange{24}{26}
\subsection*{L'équation fonctionnelle du terme à l'infini}

Soit $\bar f : \bar X \to \operatorname{Spec}\mathbb{F}_p$ une compactification,
$i : X \to \bar X$ l'ouverture, $j : Y \to \bar X$ le bord fermé. Alors
$R_\infty f(E) = R\bar f_{*}\,j^{*}Ri_{*}(E)$ et, $\bar f$ étant propre, la
dualité locale donne
\[
  D\,R_\infty f(E) = R\bar f_{*}\,Rj^{!}\bigl(Ri_{!}(DE)\bigr)
  = R_\infty f(DE)[-1],
\]
puisque $Rj^{!}Ri_{!}K \simeq j^{*}Ri_{*}K[-1]$ (triangle
$j_{*}Rj^{!} \to \mathrm{id} \to Ri_{*}i^{*}$ appliqué à $Ri_{!}K$). Avec
$DE \simeq E(n+\rho)[2n]$,
\[
  D\,R_\infty f(E) \simeq R_\infty f(E)(n + \rho)[2n - 1] .
\]
\footnote{La page écrit $[1]$ dans les deux formules encadrées, avec en marge
« $1$ ou $-1$ ! ». C'est $[-1]$, puis $[2n-1]$ après substitution. Seule la
parité du décalage compte pour la fonction $L$ --- un décalage impair
l'inverse --- et elle est la même : la suite du calcul est juste.} Le décalage
étant impair, la fonction $L$ s'inverse, et l'on obtient l'équation du bord :
\[
  L^{\infty}_E\Bigl(\frac{1}{qt}\Bigr)\,L^{\infty}_E(t)\,(-t)^{\chi_\infty(E)}\,\delta_\infty(E) = 1,
\]
avec $\chi_\infty$, $\delta_\infty$ ceux de $R_\infty f(E)$. Or
$\chi_\infty(E) = \chi(Rf_{*}E) - \chi(Rf_{!}E) = 0$ et
$\delta_\infty(E) = \delta(Rf_{*}E)\,\delta(Rf_{!}E)^{-1} = q^{\chi(E)}\delta(E)^{-2}$,
d'où
\[
  L^{\infty}_E\Bigl(\frac{1}{qt}\Bigr)L^{\infty}_E(t)\,\delta_\infty(E) = 1,
  \qquad \delta_\infty(E) = \varepsilon(E)^{-2},
\]
« qui est bien $1$ lorsque $X$ est complet ! ». Il dresse la liste des
« invariants arithmétiques intéressants » une fois connus $\chi(E)$ et $q$ :
$\delta(E) = \delta(Rf_{!}E)$, $\delta'(E) = \delta(Rf_{*}E)$,
$\varepsilon(E)$, $\delta_\infty(E)$, liés par
\[
  \begin{gathered}
  \delta(E)\,\delta'(E) = q^{\chi(E)}, \qquad
  \varepsilon(E) = \delta(E)\,q^{-\chi(E)/2} = \delta'(E)^{-1}q^{\chi(E)/2},\\
  \delta_\infty(E) = q^{\chi(E)}\delta(E)^{-2} = q^{-\chi(E)}\delta'(E)^{2} = \varepsilon(E)^{-2},
  \end{gathered}
\]
et $\varepsilon(E) = \pm 1$ si $X$ est complète.\footnote{La deuxième formule
encadrée porte $\varepsilon(E) = \delta'(E)\,q^{-\chi(E)/2}$, ce qui
contredit la première ; c'est $\varepsilon(E)^{-1}$ qui vaut
$\delta'(E)q^{-\chi(E)/2}$. Une flèche en marge, partant de ce signe $=$, dit
« variance avec $E$, $X$ !! ».}

\pagerange{28}{30}
\subsection*{Une vérification}

Posons $B(t) = A(t)\,\delta(E)(-t)^{\chi(E)}$, de sorte que
$L_E(1/qt) = B(t)L_E(t)$. Substituer $t \mapsto 1/qt$ et réinjecter donne
$B(t)B(1/qt) = 1$, et, en remplaçant $B$ par sa valeur,
$L^{\infty}_E(t)L^{\infty}_E(1/qt) = \delta(E)^2q^{-\chi(E)}$ : « c'est
précisément l'équation fonctionnelle explicitée plus haut ».\footnote{Dans la
ligne intermédiaire, le facteur $(1/qt)^{\chi(E)}$ doit porter le signe
$(-1)^{\chi(E)}$ ; il se compense avec celui de $(-t)^{\chi(E)}$, et le
résultat est juste.} Avec $\xi_E(t) = t^{\chi(E)/2}L_E(t)$,
\[
  \xi_E\Bigl(\frac{1}{qt}\Bigr) = A(t)\,\varepsilon'(E)\,\xi_E(t),
  \qquad A(t) = L^{\infty}_E\Bigl(\frac{1}{qt}\Bigr)^{-1},
  \qquad \varepsilon'(E) = (-1)^{\chi(E)}\varepsilon(E),
\]
et si $X$ est complète, $A = 1$ et $\varepsilon'(E) = \pm 1$.\footnote{La page 30
écrit $q = p^{n\rho}$ ; c'est $q = p^{n+\rho}$, comme aux pages 26 et 28. Un
petit exposant sur $A$ ne se lit pas.}

\pagerange{32}{36}
\section*{5. Une courbe, un faisceau au point générique (pages 32 à 36)}

Soit $X$ une courbe projective, lisse, géométriquement irréductible sur
$\mathbb{F}_p$, de corps de fonctions $K$, $\eta = \operatorname{Spec}K$, et
$E_\eta$ un $\mathbb{Q}_\ell$-faisceau sur $\eta$ --- une représentation
$\ell$-adique de $\operatorname{Gal}(K^{s}/K)$. On veut lui associer une
fonction $L$, ce qui « est raisonnable, car $K$ détermine $X$ ». Avec
$i : \eta \to X$, posons $L^{*}_{E_\eta} = L_{i_{*}E_\eta}$, la fonction $L$
usuelle, dont le facteur en $x$ est pris sur les invariants de l'inertie.
On a $D(i_{*}E_\eta) = i_{*}(\check E_\eta)(1)[2]$,\footnote{La page omet le
décalage $[2]$ ; il est pair et sans effet sur les fonctions $L$.} et (1)
donne
\[
  L^{*}_{\check E_\eta}\Bigl(\frac{1}{pt}\Bigr)
  = (-t)^{\chi^{*}(E_\eta)}\,\delta^{*}(E_\eta)\,L^{*}_{E_\eta}(t),
  \qquad \chi^{*} = \chi(i_{*}E_\eta),\ \delta^{*} = \delta(i_{*}E_\eta).
\]
\footnote{La page écrit l'exposant $\chi^{*}(\check E_\eta)$ ; il est égal à
$\chi^{*}(E_\eta)$, puisque $D$ conserve la caractéristique d'Euler.} Si
$\check E_\eta \simeq E_\eta(\rho)$, alors, avec $q = p^{1+\rho}$,
\[
  L^{*}_{E_\eta}\Bigl(\frac{1}{qt}\Bigr)
  = (-t)^{\chi^{*}}\delta^{*}\,L^{*}_{E_\eta}(t), \qquad
  \xi^{*}\Bigl(\frac{1}{qt}\Bigr) = \varepsilon'^{*}\,\xi^{*}(t), \qquad
  \varepsilon'^{*} = \delta^{*}q^{-\chi^{*}/2}(-1)^{\chi^{*}} .
\]

\emph{Une remarque} : $L^{*}$ n'est pas additive. D'une suite exacte
$0 \to E'_\eta \to E_\eta \to E''_\eta \to 0$ on tire
\[
  0 \to i_{*}E'_\eta \to i_{*}E_\eta \to i_{*}E''_\eta \to Q \to 0,
  \qquad Q \simeq \bigoplus_x j_{x*}Q_x, \qquad
  Q_x = \operatorname{Coker}\bigl(E_\eta^{I_x} \to E_\eta''^{I_x}\bigr),
\]
somme sur les points de ramification, d'où
\[
  L^{*}_{E_\eta}(t) = L^{*}_{E'_\eta}(t)\,L^{*}_{E''_\eta}(t)\prod_x P_{Q_x}(t),
  \qquad P_{Q_x}(t) = \det\bigl(1 - F_x t^{\deg x} \mid Q_x\bigr).
\]
\footnote{La page écrit $P_{Q_x}(t)$ sans préciser le degré du point ; le
facteur est en $t^{\deg x}$. Une marge, en partie illisible, rappelle la dualité
entre $H^2(\bar X, E)$ et $H^0$ du dual tordu.}

On peut songer à la rendre additive en passant aux facteurs de composition
--- définir $L$ sur les $E_\eta$ simples et prolonger par linéarité, ou prendre
pour facteur local en $x$ la somme des invariants des facteurs d'une suite de
composition. « Ceci n'est pas encore satisfaisant » : ces facteurs ne sont pas
de nature locale en $(E, x)$, car la suite de composition de $E_\eta$ comme
représentation globale n'a rien à voir avec celle de sa restriction au groupe
de décomposition en $x$, et ne commute pas à l'extension étale.\footnote{La
fin de la phrase précise « (i.e.\ ext.\ résiduellement … triviale) », avec un
mot ajouté illisible.}

\pagerange{38}{46}
\section*{6. La moyenne sur l'inertie (pages 38 à 46)}

La page 38 s'intitule « Suite de la remarque ». Soit $D$ un groupe, $I$ un
sous-groupe distingué, $E$ une représentation de dimension finie de $D$ sur un
corps $k$, et $I_1 \subset I$ un sous-groupe d'indice fini, distingué dans
$D$, qui opère sur $E$ de façon \emph{unipotente}.\footnote{La page note $U$
ce sous-groupe et $V \subset U$ le suivant ; on écrit $I_1$, $I_2$, parce que
$U$ servira pour un ouvert de la courbe. « Invariant par $D$ » est ajouté
au-dessus de la ligne. La marge observe que, dans le cas galoisien, un $I_1$
ouvert contient un sous-groupe ouvert invariant par $D$. C'est vrai, mais
c'est le théorème de monodromie qui le donne et non la théorie des groupes
profinis : l'ensemble des éléments de $I$ qui opèrent de façon unipotente est
le noyau de la partie d'image finie de la représentation, ouvert et stable par
$D$.} La suite $E^{I_1} \subset \cdots$, de quotients successifs
$E^{I_1}, (E/E^{I_1})^{I_1}, \dots$, est une suite de composition stable par
$D$, sur les quotients de laquelle $I_1$ opère trivialement. Dans le groupe de
Grothendieck, $E$ s'identifie donc à une représentation de $D/I_1$.

Sur $\mathrm{Mod}(D/I_1, k)$, la moyenne
\[
  E \longmapsto E^{I} = \operatorname{Im}\pi, \qquad
  \pi = \frac{1}{N}\sum_{g \in I/I_1} g, \qquad N = [I : I_1],
\]
est un foncteur exact vers $\mathrm{Mod}(D/I, k)$ dès que $N$ est inversible
dans $k$, en particulier en caractéristique $0$.\footnote{La page dit « c'est
le \emph{seul} foncteur exact », lecture incertaine ; on ne retient que
l'exactitude.} Il induit
$\natural_{I_1} : R(D/I_1, k) \to R(D/I, k)$, compatible, pour $I_2 \subset
I_1$, à l'inflation $R(D/I_1, k) \to R(D/I_2, k)$. On obtient ainsi une
opération additive $\operatorname{cl}(E) \mapsto \operatorname{cl}(E)^{\natural}$,
définie dès qu'un sous-groupe d'indice fini de $I$ opère de façon unipotente,
compatible aux représentations continues et à l'extension du corps de
base.\footnote{Le symbole de l'opération a la forme d'un éclair ; la
transcription le rend par $\natural$ et l'on garde ce signe. Un paragraphe
biffé de la page 42 commençait le cas $D/I \simeq \mathbb{Z}$.} Sa trace se
calcule sans quotienter : pour $\varphi \in D/I$,
\[
  \operatorname{Tr}\bigl(\varphi \mid E^{\natural}\bigr)
  = \frac{1}{N}\sum_{\substack{\psi \in D/I_1 \\ \psi \mapsto \varphi}}
    \operatorname{Tr}(\psi \mid E),
\]
où $\operatorname{Tr}(\psi \mid E)$, pour $\psi \in D/I_1$, est la trace d'un
relèvement quelconque de $\psi$ dans $D$ --- elle n'en dépend pas, puisque
$I_1$ opère de façon unipotente. La formule, vraie quand $I_1$ opère
trivialement, s'étend parce que ses deux membres sont additifs.

Revenons au cas géométrique : $D = D_x$ est le groupe de décomposition en un
point fermé $x$, $I = I_x$ l'inertie, $D_x/I_x \simeq \hat{\mathbb{Z}}$,
engendré par le Frobenius $F_x$. L'hypothèse --- un sous-groupe ouvert de
$I_x$ opère de façon unipotente --- est toujours satisfaite : c'est le théorème
de monodromie $\ell$-adique de Grothendieck.\footnote{Le dossier ne le cite pas
et le suppose implicitement en posant l'opération sur tout $E_\eta$. Le
théorème figure dans l'appendice de Serre--Tate, \emph{Good reduction of abelian
varieties} (1968) ; on ne sait pas, faute de date, s'il est antérieur ou
postérieur à ces pages.} En termes d'aujourd'hui : si $E$ correspond à la
représentation de Weil--Deligne $(r, N)$, $E^{\natural(x)}$ est $(E, r)^{I_x}$,
les invariants de l'inertie quand on oublie l'opérateur de monodromie $N$, alors
que la fibre de $i_{*}E_\eta$ en $x$ est $E^{I_x} = (\ker N)^{I_x}$.\footnote{Ce
langage est de Deligne (1973), et n'est pas sur la page ; la traduction est
nôtre.} On prend $E_\eta^{\natural(x)}$ comme fibre modifiée en $x$, et l'on
définit le faisceau virtuel sur $X$
\[
  E_\eta^{\natural} = u_{!}(E^{\circ}) + \sum_{x \in Y} j_{x*}\bigl(E_\eta^{\natural(x)}\bigr),
\]
où $X^{\circ} \subset X$ est un ouvert dense sur lequel $E_\eta$ est non ramifié,
$E^{\circ}$ le faisceau lisse qu'il y définit, $u : X^{\circ} \to X$,
$Y = X - X^{\circ}$ et $j_x : \operatorname{Spec}k(x) \to X$.\footnote{La page
note $U$, $E_U$, $i : U \to X$ ; et elle écrit
« $E_\eta(x) = E_\eta^{\natural(x)}$ », alors que la page 50 réserve
$E_\eta(x)$ à la fibre ordinaire $E^{I_x}$. On suit la page 50.} On note
$L^{\natural}_{E_\eta} = L_{E_\eta^{\natural}}$,
$\chi^{\natural}(E_\eta) = \chi(E_\eta^{\natural})$,
$\delta^{\natural}(E_\eta) = \delta(E_\eta^{\natural})$. Par construction,
$E_\eta \mapsto E_\eta^{\natural}$ est additif et local en $(E, x)$ : c'est ce
qui manquait à la page 36.

\pagerange{48}{58}
\section*{7. La dualité de $E_\eta^{\natural}$ (pages 48 à 58)}

\subsection*{Le calcul de la page}

Comme la moyenne locale commute au dual ($(\check E)^{\natural(x)} =
(E^{\natural(x)})^{\vee}$, l'inertie opérant par un groupe fini) et que
$D(E^{\circ}) = \check E^{\circ}(1)[2]$, on a dans le groupe de
Grothendieck\footnote{La page omet les décalages $[2]$, pairs ; et la lecture
« $\mathbb{Q}_\ell$ » dans $\mu_x$ est douteuse, mais c'est la seule qui rende
la formule juste.}
\[
  D(E_\eta^{\natural}) = Ru_{*}(\check E^{\circ})(1) + \sum_x j_{x*}\bigl(\check E_\eta^{\natural(x)}\bigr)
  = \check E_\eta^{\natural}(1) + \sum_{x \in X^{(0)}} j_{x*}\bigl(\mu_x(E)\bigr),
\]
\[
  \mu_x(E) = \check E_\eta^{\natural(x)} \otimes \bigl(\mathbb{Q}_\ell - \mathbb{Q}_\ell(1)\bigr)
           + j_x^{*}Ru_{*}(\check E^{\circ})(1).
\]
Par dualité locale, $j_x^{*}Ru_{*}(\check E^{\circ}) = \check E_\eta(x) -
E_\eta(x)^{\vee}(-1)$, où $E_\eta(x) = E^{I_x}$. Posant
\[
  \alpha_x(E_\eta) = E_\eta^{\natural(x)} - E_\eta(x),
  \qquad\text{il vient}\qquad
  \mu_x(E) = \alpha_x(E_\eta)^{\vee} - \alpha_x(\check E_\eta)(1),
\]
et $\mu_x(\check E) = -\mu_x(E)^{\vee}(1)$. « NB. $\alpha_x$ n'est \emph{pas}
additif en $E_\eta$, \emph{mais} $\mu_x$ \emph{l'est}. » Si
$\check E_\eta \simeq E_\eta(\rho)$,
$\mu_x(E) = \alpha_x(E_\eta)^{\vee} - \alpha_x(E_\eta)(\rho + 1)$.\footnote{Ici
et aux pages 52 à 56, la lettre ressemble à un $\varphi$ ; c'est le $\rho$
des pages précédentes.}

\subsection*{Ce que la page ne dit pas : $\mu_x$ est nul}

En fait $\mu_x(E) = 0$ pour tout $E_\eta$ --- c'est notre observation, la page
ne la fait pas. Pour une représentation $G$ de $D_x$ sur laquelle un
sous-groupe ouvert $I_1$ de $I_x$ opère de façon unipotente, on a dans le groupe
de Grothendieck des $F_x$-modules
\[
  \bigl[H^0(I_x, G)\bigr] - \bigl[H^1(I_x, G)\bigr] = \bigl[G^{\natural}\bigr] - \bigl[G^{\natural}(-1)\bigr] .
\]
Les deux membres sont additifs en $G$ (à gauche parce que
$\operatorname{cd}_\ell(I_x) = 1$) ; il suffit donc de le voir quand $I_1$ opère
trivialement. Alors $H^0 = G^{I_x} = G^{\natural}$ et
$H^1(I_x, G) = \operatorname{Hom}(I_1, G)^{I_x/I_1} = G(-1)^{I_x/I_1} =
G^{\natural}(-1)$, puisque le plus grand quotient pro-$\ell$ abélien de $I_1$ est
$\mathbb{Z}_\ell(1)$. Comme $j_x^{*}Ru_{*}(\check E^{\circ})$ est exactement
$[H^0(I_x, \check E)] - [H^1(I_x, \check E)]$, on obtient
$j_x^{*}Ru_{*}(\check E^{\circ}) = \check E^{\natural(x)} - \check E^{\natural(x)}(-1)$,
et
\[
  \mu_x(E) = \check E^{\natural(x)} - \check E^{\natural(x)}(1)
           + \check E^{\natural(x)}(1) - \check E^{\natural(x)} = 0 .
\]
Autrement dit $\alpha_x(\check E_\eta) = \alpha_x(E_\eta)^{\vee}(-1)$, et
\[
  D(E_\eta^{\natural}) = \check E_\eta^{\natural}(1)[2]
\]
dans le groupe de Grothendieck : la fonction $L^{\natural}$ a la même dualité
que $L^{*}$, sans terme correctif, et elle est de plus additive et
locale.\footnote{Ce qui rend la chose plausible : la dualité de Poincaré pour
$i_{*}E_\eta$ est exacte, $D(i_{*}E_\eta) = i_{*}\check E_\eta(1)[2]$, et
$E_\eta^{\natural} - i_{*}E_\eta = \sum_x j_{x*}\alpha_x(E_\eta)$, dont le dual
est $\sum_x j_{x*}\alpha_x(E_\eta)^{\vee} = \sum_x j_{x*}\alpha_x(\check E_\eta)(1)$.
En termes de Weil--Deligne, $\alpha_x(E)$ est la classe de
$(E/\ker N)^{I_x}$, que $N$ identifie à $(\operatorname{Im}N)^{I_x}(-1)$ ;
c'est l'orthogonalité de $\ker N$ et de $\operatorname{Im}N^{\vee}$ qui fait
le reste.}

\pagerange{52}{58}
\subsection*{L'équation fonctionnelle, et la question de la page 52}

Avec les $\mu_x$ de la page, l'équation (1) appliquée à $E_\eta^{\natural}$
s'écrit, en posant $\lambda_x(t) = L_{j_{x*}\mu_x(E)}(t)$,
\[
  L^{\natural}_{E_\eta}\Bigl(\frac{1}{pt}\Bigr)
  = \prod_x \lambda_x(pt)\,\bigl(-pt\bigr)^{\chi^{\natural}(E_\eta)}
    \delta^{\natural}(E_\eta)^{-1}\,L^{\natural}_{\check E_\eta}(t),
\]
puis, si $\check E_\eta \simeq E_\eta(\rho)$ et $q = p^{1+\rho}$,
\[
  L^{\natural}_{E_\eta}\Bigl(\frac{1}{qt}\Bigr)
  = \Bigl(\prod_x \lambda_x(qt)\Bigr)\,\delta^{\natural}(E_\eta)^{-1}
    \bigl(-qt\bigr)^{\chi^{\natural}(E_\eta)}\,L^{\natural}_{E_\eta}(t) .
\]
\footnote{La page écrit $\prod_x \lambda_x(E)$, puis $\prod_x \lambda_x(E)(t)$,
sans l'argument $pt$ puis $qt$ que les substitutions $t \mapsto pt$ et
$t \mapsto p^{\rho}t$ imposent.} Elle développe chaque $\lambda_x$ en un
monôme $(-t)^{-\chi_x}\delta_x$ fois un quotient de fonctions $L$ de
$\alpha_x$, où $\chi_x$ et $\delta_x$ sont la caractéristique d'Euler et le
déterminant de $j_{x*}\alpha_x(E_\eta)$ vu sur
$\operatorname{Spec}\mathbb{F}_p$ ; la page 54 regroupe le tout sous la forme
$L(1/qt) = A(t)L(t)$, et ajoute : « Ayant écrit $\lambda_x$ comme un monôme, on
en perd le \emph{caractère additif}, il vaut sans doute mieux ne pas faire ce
changement ».\footnote{La lettre notée $\chi_x$ est surchargée aux pages 52 et
56 et pourrait être un $\beta_x$. On prend $\chi_x$ et $\delta_x$ après
image directe sur $\operatorname{Spec}\mathbb{F}_p$, ce qui absorbe le signe
$(-1)^{(\deg x - 1)\operatorname{rg}}$ que le déterminant d'une induite
introduit ; la page ne le dit pas.} Puis, entre crochets, une question :
\[
  \delta^{\natural}(E_\eta)^2 \overset{?}{=} q^{\chi^{\natural}(E_\eta)} \quad ??
\]

La page 56 la ramène à des termes locaux. Lisant les déterminants sur la
formule de dualité, avec $\delta(\mu_x) = \delta_x^{-2}q^{\chi_x}$, il trouve
\[
  \varepsilon^{\natural}(E_\eta)^2 = \frac{\delta^{\natural}(E_\eta)^2}{q^{\chi^{\natural}(E_\eta)}}
  = \prod_x \frac{\delta_x(E_\eta)^2}{q^{\chi_x(E_\eta)}},
  \qquad
  \varepsilon^{\natural}(E_\eta) = \pm \prod_x \delta_x(E_\eta)\,q^{-\chi_x(E_\eta)/2} .
\]
\footnote{La page ajoute une variante,
$\varepsilon^{\natural} = \pm\bigl(\prod_x \delta_x p^{-\rho\chi_x/2}\bigr)p^{-\chi/2}$
avec $\chi = \sum_x \chi_x$, qui est la même chose.} Comme $\mu_x = 0$, chaque
facteur vaut $1$, $\lambda_x = 1$, et la réponse à la question est
\emph{oui} : pour $E_\eta$ autodual de poids $\rho$,
$\delta^{\natural}(E_\eta)^2 = q^{\chi^{\natural}(E_\eta)}$, et
\[
  L^{\natural}_{E_\eta}\Bigl(\frac{1}{qt}\Bigr)
  = \varepsilon^{\natural}(E_\eta)^{-1}\bigl(-q^{1/2}t\bigr)^{\chi^{\natural}(E_\eta)}
    L^{\natural}_{E_\eta}(t), \qquad \varepsilon^{\natural}(E_\eta) = \pm 1 .
\]

Restent les problèmes de la page 54 :
\begin{itemize}
  \item[a)] dans les cas géométriques classiques, construire les
    $\alpha_x(E_\eta)$, et en déterminer $\chi_x$, $\delta_x$ et les valeurs
    absolues des valeurs propres du Frobenius ;
  \item[b)] une formule générale « (Chafarevitch--)Ogg » pour
    $\chi^{\natural}(E_\eta)$ ;
  \item[c)] une formule pour $\delta^{\natural}(E_\eta)$.
\end{itemize}
Les flèches de la marge rangent $\chi_x$ et $\chi^{\natural}$ parmi les
invariants géométriques, $\delta_x$ et $\delta^{\natural}$ parmi les
arithmétiques. Page 58 : « c) est ramené (\emph{moyennant un signe}) au cas de
b), donc à déterminer $\chi^{\natural}(E_\eta)$. Or cette détermination est
faite, en utilisant les représentations d'Artin … » --- et la page s'arrête.
C'est exactement ce que donne $\delta^{\natural\,2} = q^{\chi^{\natural}}$.
La formule visée en b) est celle de Grothendieck--Ogg--Chafarevitch :
$\chi(\bar X, i_{*}E_\eta) = (2 - 2g)\operatorname{rg}E_\eta - \sum_x \deg x
\bigl(\operatorname{rg}E_\eta - \operatorname{rg}E_\eta^{I_x} +
\operatorname{Sw}_x(E_\eta)\bigr)$, d'où
$\chi^{\natural} = \chi(\bar X, i_{*}E_\eta) + \sum_x \deg x\,\operatorname{rg}\alpha_x$.
\footnote{La lecture « Chafarevitch » est incertaine. Ogg (1962) et
Chafarevitch (1961) l'ont pour les courbes elliptiques et les espaces
principaux homogènes ; Grothendieck l'établit pour les faisceaux $\ell$-adiques
sur une courbe (exposé Bourbaki de Raynaud n° 286, 1965 ; SGA~5, exposé~X), au
moyen des représentations de Swan et d'Artin --- ce que la page 58 appelle
« les représentations d'Artin ». Le signe qui reste, lui, est le signe de
l'équation fonctionnelle, que la théorie des constantes locales de Langlands et
Deligne (1972--1973) et la formule du produit de Laumon (1987) écrivent comme
produit de facteurs locaux ; le dossier ne l'aborde pas.}

\pagerange{60}{60}
\section*{8. Cohomologie à l'infini (page 60)}

Retour aux pages 18 à 26. Soit $f$ compactifiable en
$\bar f : \bar X \to \operatorname{Spec}\mathbb{F}_p$, $i : X \to \bar X$,
$j : Y = \bar X - X \to \bar X$. Le triangle
$Ri_{!}F \to Ri_{*}F \to j_{*}\bigl((Ri_{*}F)|_Y\bigr) \xrightarrow{+1}$ donne,
par $R\bar f_{*}$, la suite exacte longue
\[
  \cdots \to R^i f_{!}(F) \to R^i f_{*}(F) \to R^i_\infty f(F) \to R^{i+1}f_{!}(F) \to \cdots,
  \qquad R^i_\infty f(F) = R^i\bar f_{*}\bigl((Ri_{*}F)|_Y\bigr).
\]
Il considère aussi le triangle
$j_{*}Rj^{!}(Ri_{!}F) \to Ri_{!}F \to Ri_{*}F \xrightarrow{+1}$, et veut
identifier $(Ri_{*}F)|_Y$, « à décalage près », à $Rj^{!}Ri_{!}F$. Les deux
triangles le donnent : $Rj^{!}Ri_{!}F \simeq (Ri_{*}F)|_Y[-1]$, l'isomorphisme
dont on s'est servi pour la dualité de la page 24.\footnote{La page s'arrête sur
un tiret, et ses derniers verbes sont d'une lecture incertaine ; on donne
l'identification qu'elle annonce, avec le décalage qu'elle ne précise pas.}

\pagerange{62}{68}
\section*{9. Simplifier une équation fonctionnelle (pages 62 à 68)}

Soit $L(t)$ une fonction rationnelle vérifiant $L(1/qt) = A(t)L(t)$, avec $A$
« de niveau inférieur ».\footnote{L'exposant de $p$ dans $q$ est noyé sous une
tache d'encre ; la page 66 substitue $1/(p^{n+\rho}t)$ là où la page 68 écrit
$1/qt$, donc $q = p^{n+\rho}$. « Niveau » est son mot, entre guillemets, et
n'est pas défini ; on en donne plus bas le sens que ces pages lui font jouer.}
On cherche $\lambda$ tel que $L' = \lambda L$ vérifie $L'(1/qt) = c\,L'(t)$,
c'est-à-dire
\[
  \frac{\lambda(1/qt)}{\lambda(t)} = c\,A(t)^{-1} .
\]
Deux solutions diffèrent par un $\alpha$ tel que $\alpha(1/qt) = a\,\alpha(t)$, et
$a = \pm 1$ nécessairement, puisque $t \mapsto 1/qt$ est une involution.
Faisant opérer $\mathbb{Z}/2\mathbb{Z}$ sur les fractions rationnelles non nulles
par cette involution, la question devient celle d'écrire un $1$-cocycle comme
un cobord, $\lambda^{g}\lambda^{-1} = c\,A^{-1}$. Il le peut toujours :
$A(t)A(1/qt) = 1$ (substituer deux fois), donc $c A^{-1}$ est un cocycle pour
$c = \pm 1$, et le théorème 90 de Hilbert pour l'extension quadratique
$\mathbb{Q}(t)/\mathbb{Q}(t)^{\mathbb{Z}/2}$ dit qu'il est un cobord. La page
donne d'ailleurs la solution triviale, $\lambda = 1/L$, $L' = 1$.\footnote{Le
nom du théorème 90 n'est pas sur la page ; l'action de $\mathbb{Z}/2\mathbb{Z}$
l'est, et un encadré « $\lambda^{g}\lambda^{-1} = $ » laissé sans second
membre.} Toute la question est donc d'imposer à $\lambda$ d'être de « niveau »
strictement inférieur à celui de $L$.

Pour préciser, dans un passage ensuite barré de trois traits, il suppose
$L(0) = 1$ et les zéros et pôles de module $p^{-i/2}$, $\rho \le i \le \rho + 2n$,
et décompose
\[
  L(t) = L_\rho(t)\cdots L_{\rho+n}(t)\;
         L'_{\rho+n-1}(pt)\,L'_{\rho+n-2}(p^2t)\cdots L'_\rho(p^nt),
\]
où les valeurs propres inverses de $L_i$ et de $L'_i$ sont des entiers
algébriques de module $p^{i/2}$ ; celles de $L'_{\rho+n-k}(p^kt)$ sont donc de
module $p^{(\rho+n+k)/2}$. Le « niveau » d'un facteur est ainsi le poids qu'il
a avant torsion de Tate : $L$ est de niveau $\le \rho + n$ et de poids entre
$\rho$ et $\rho + 2n$.\footnote{Interprétation nôtre, qui rend compte de
toutes les occurrences du mot. Le passage barré écrivait « de niveau
$\le \rho + n + 1$ » pour $\lambda$ ; la page 66, non barrée, dit
$\le \rho + n - 1$, et c'est elle qu'on suit. Les arguments $p^2t$, $p^nt$ se
lisent mal.} Page 66, il s'autorise à corriger par des $\lambda$ de même forme
mais de niveau $\le \rho + n - 1$, suppose que $A$ est lui aussi de ce type et
que chaque $L_i$, $L'_i$ vérifie l'équation fonctionnelle de son poids,
$L_i(1/p^it) = \delta_i t^{\chi_i}L_i(t)$ et de même pour $L'_i$. Alors
\[
  \frac{L(1/qt)}{L(t)} = \delta\,t^{\chi}
  \prod_{i<\rho+n}\frac{L_i(p^{\rho+n-i}t)}{L_i(t)}
  \prod_{i}\frac{L'_i(t)}{L'_i(p^{\rho+n-i}t)} = A(t),
\]
et le facteur de niveau maximal, $L_{\rho+n}$, a disparu de $A$ : c'est le sens
précis de « $A$ de niveau inférieur ».\footnote{La page 66 omet les facteurs
$p^{(\rho+n-i)\chi_i}$ que la substitution produit ; ils sont absorbés dans la
constante $\delta$. Elle écrit aussi $t^{\chi_\rho}$ devant $L_{\rho+n}$, où
l'on attend $\chi_{\rho+n}$, et une suite d'indices pour $\lambda$ qu'on ne
reproduit pas, la page l'ayant écrite de façon incohérente.}

Le cas $n = 1$ est commencé page 68 : $L = L_\rho(t)\,L_{\rho+1}(t)\,L'_\rho(pt)$,
avec les facteurs $L_\rho(t)$ et $L'_\rho(pt)$ soulignés d'une accolade,
puis réécrits seuls --- et la page s'arrête. La conclusion vers laquelle elle
va est immédiate : $\lambda = \bigl(L_\rho(t)L'_\rho(pt)\bigr)^{-1}$ est de
niveau $\rho = \rho + n - 1$, et $L' = \lambda L = L_{\rho+1}$ vérifie une
équation fonctionnelle sans facteur, de même $q$ ; en général,
$\lambda = \bigl(\prod_{i<\rho+n}L_i(t)\prod_i L'_i(p^{\rho+n-i}t)\bigr)^{-1}$
laisse $L' = L_{\rho+n}$.\footnote{Cette conclusion est nôtre ; la page ne
l'écrit pas.} Ce que le dossier cherchait, à travers tout cela, c'est une
fonction $L$ qui n'ait gardé que sa partie de plus haut niveau --- celle qui
porte, à elle seule, une équation fonctionnelle exacte.

\end{document}
